% 第 16 章 （导出）几何 Langlands 中的缺陷
% 原文：Chapter 16: (Derived) defects in Geometric Langlands
%        （pp. 265–279，印刷页码 264–278）
% 第 35 讲；2021 年 8 月 20 日
% ============================================================================
\chapter{（导出）几何 Langlands 中的缺陷}\label{ch:16}

\begin{flushright}\small 第 35 讲；2021 年 8 月 20 日\end{flushright}

\begin{keytakeaways}
本章讨论几何 Langlands 中的“导出缺陷”。先从几何 Satake 出发：线算子在 $A$-侧由 $\Shv(LG_+\backslash LG/LG_+)$ 给出，在 $B$-侧由 $\Rep(G^\vee)$ 给出；两者在导出情形下都升级为范畴。四维理论 $\cZ$ 的局部算子 $\cZ(S^3)$ 构成 $\mathbb{E}_4$-代数，线算子 $\cZ(S^2)$ 构成 $\mathbb{E}_3$-范畴，于是得到 $B$-侧的真空模空间 $\cM^1_{B_{G^\vee}}=\fg^{\vee*}[2]/G^\vee$ 及其“幂零锥”$\cM^0=\mathfrak{h}^{\vee*}[2]/W$。核心工具是 Koszul 对偶与导出自交：$\Loc_{G^\vee}(S^2)=\Omega G^\vee/G^\vee$，其上的ind-凝聚层由 $S=\mathrm{Sym}(V[-2])$ 的模刻画。由此得到 Arinkin--Gaitsgory 的奇异支撑理论以及 Beraldo 的定理：由线算子构造的可观测量正是 $T^*[1]\Loc_{G^\vee}$ 的形变量子化模。本章最后讨论余维一缺陷（畴壁与边界条件）、Langlands 函子性，以及“边界条件 $\leftrightarrow$ 移位哈密顿 $G^\vee$-空间”这一对应。
\end{keytakeaways}

\section{通过辛几何研究 TFT}
\begin{sectionkey}
本节从几何 Satake 出发：它把 $\mathcal{A}_G$ 的线算子等同于 $\mathcal{B}_{G^\vee}$ 的线算子，后者由 $\Rep(G^\vee)=\QC(\mathrm{pt}/G^\vee)$ 给出，于是（非导出）真空模空间是 $\cM^1=\mathrm{pt}/G^\vee$。由于单位线算子的自同态平凡，非导出情形下看不到局部算子；导出版本把线算子升级为 $\mathbb{E}_3$-范畴，并把真空模空间变为移位辛叠 $\fg^{\vee*}[2]/G^\vee=T^*[3](\mathrm{pt}/G^\vee)$。其来源是 $\Loc_{G^\vee}(S^2)$ 作为 $G^\vee$ 中一点的自交，即基于一点的环路空间 $\Omega G^\vee/G^\vee$，其上的层由 Koszul 对偶描述。
\end{sectionkey}

回想几何 Satake 把理论 $\mathcal{A}_G$ 中的线算子等同于理论 $\mathcal{B}_{G^\vee}$ 中的线算子。也就是说，它给出等同
\begin{equation}\tag{16.1}
\begin{array}{c}
\mathcal{A}_G(S^2)=\Shv(LG_+\backslash LG/LG_+)\\[3pt]
\big\Downarrow\\[3pt]
\mathcal{B}_{G^\vee}(S^2)=\QC(\Loc_{G^\vee}S^2).
\end{array}
\end{equation}
$B$-侧可以表达为：
\begin{equation}\tag{16.2}
\mathcal{B}_{G^\vee}(S^2)=\QC(\Loc_{G^\vee}S^2)=\QC(\mathrm{pt}/G^\vee)=\Rep(G^\vee).
\end{equation}
注意 $A$-侧 $\Shv(LG_+\backslash LG/LG_+)$ 上的乘法由卷积给出，而 $B$-侧 $\Rep(G^\vee)$ 上的乘法是张量积。

理解这个结果的一种方式是它断言：
\begin{equation}\tag{16.3}
\cM^1_B=\cM^1_A=\mathrm{pt}/G^\vee.
\end{equation}
我们还可以问这些理论中的局部算子。想法是：理论 $\cZ$ 的局部算子由 $\cZ(S^3)$ 构成；等价地，它们是单位线算子的 $\End$，即 $\cZ(S^2)=\Rep(G^\vee)$ 中平凡表示的 $\End$。其想法是，对单位取 $\End$ 是悬挂（suspension）的一种版本。但单位没有自同态：
\begin{equation}\tag{16.4}
\End_{\Rep(G^\vee)}(\mathrm{triv})=k.
\end{equation}
这是我们在几何 Langlands 中谈局部算子谈得不多的一大原因。

\begin{srcrem}{113}
今天我们不会多谈面缺陷，但它们正是分歧数据。
\end{srcrem}

\subsection{导出版本}

从物理的观点看，导出版本实际上更自然。下面是 TFT 中的一个一般法则。设 $\cZ$ 是四维 TFT。则局部算子 $\cZ(S^3)$ 构成 $\mathbb{E}_4$-代数，即某种类似奇 Poisson 代数的东西。所以它略微非交换。于是
\begin{equation}\tag{16.5}
\cM^0=\Spec\cZ(S^3)
\end{equation}
是奇 Poisson 簇。类似地，线算子构成 $\mathbb{E}_3$-代数：
\begin{equation}\tag{16.6}
\cZ(S^2),
\end{equation}
因为它附着在 $2$-球面上。不过现在它是一个范畴而非代数。所以它是一个带有并非完全交换的乘法的范畴。在阿贝尔层面，$\mathbb{E}_3$ 与对称幺半范畴相同。因此，若不进入导出情形，就“感觉不到”这种非交换性。$\mathbb{E}_2$ 无需导出就能被探测到（作为编织），但 $\mathbb{E}_3$ 不行。这个 $\mathbb{E}_4$-代数与 $\mathbb{E}_3$-范畴之间的关系是：
\begin{equation}\tag{16.7}
\cZ(S^3)=\End 1_{\cZ(S^2)}.
\end{equation}
这就是 $\mathbb{E}_n$ 与 $\mathbb{E}_m$ 之间的可加关系。% [Tutor Note: “可加关系”指把链环维数降低一维时，算子代数降一阶：$\mathbb{E}_4$-代数 $\cZ(S^3)$ 正是 $\mathbb{E}_3$-范畴 $\cZ(S^2)$ 中单位对象的自同态代数。注意 (16.6) 与 (16.7) 中 $\cZ(S^2)$ 先是范畴、后又作为对象出现，这是笔记的简写；严格写法是对范畴取单位对象 $1_{\cZ(S^2)}$。]

\begin{srcrem}{114}
我们可以照此继续，例如考虑面缺陷，它们构成 $\mathbb{E}_2$ $2$-范畴，但这里不做。
\end{srcrem}

第 $0$ 个模空间是：
\begin{equation}\tag{16.8}
\cM^0=\Spec\cZ(S^3),
\end{equation}
而由 Tannakian 重构得到的（更精细的）版本是：
\begin{equation}\tag{16.9}
\cZ^1=\text{“}\Spec\text{”}\cZ(S^2).
\end{equation}
回想 Tannakian 重构取一个张量范畴 $\cC$，构造 $\Spec\cC$，其定义为：
\begin{equation}\tag{16.10}
(\Spec\cC)(R)=\Hom_{\otimes}(\cC,R\text{-}\mathrm{mod}).
\end{equation}
现在 $R$ 是 $\mathbb{E}_4$-代数，而 $\cC$ 是 $\mathbb{E}_3$-范畴，所以谱是：
\begin{equation}\tag{16.11}
(\Spec\cC)(R)=\Hom_{\mathbb{E}_3}(\cC,R\text{-}\mathrm{mod}).
\end{equation}
也就是说，这在导出世界里说得通。$\cM^0$ 是仿射 $\mathbb{E}_4$-概形，而 $\cM^1$ 是（非仿射的）$\mathbb{E}_4$ 叠/概形。它们之间的关系是：
\begin{equation}\tag{16.12}
\cM^0=\Spec\End(1_{\cC})=\Spec\End(\cO_{\cM^1})=\Spec\Gamma(\cO_{\cM^1})=\mathrm{Aff}(\cM^1).
\end{equation}

\subsection{几何 Langlands 中的移位辛空间}

我们将聚焦 $B$-侧，因为“答案”（来自 $A$-侧）在那里。问题是：与理论 $\mathcal{B}_{G^\vee}$ 相联系的真空模空间是什么？也就是说，我们想找出：
\begin{equation}\tag{16.13}
\cM^1_{B_{G^\vee}}\to\cM^0_{B_{G^\vee}}.
\end{equation}
在非导出版本中这是：
\begin{equation}\tag{16.14}
\mathrm{pt}/G^\vee\to\mathrm{pt}.
\end{equation}
现在我们要做导出版本。

\begin{srcthm}{16}
$\cM^1_{B_{G^\vee}}=\fg^{\vee*}[2]/G^\vee$。
\end{srcthm}

也就是说，结构层是
\begin{equation}\tag{16.15}
\cO_{\cM^1_B}=\mathrm{Sym}\,\fg^\vee[-2]
\end{equation}
作为 $G^\vee$ 的表示。暂且不管移位 $2$，这在几何上是非常熟悉的东西：它是带有 $G^\vee$ 作用的余伴随表示，是表示论的基本对象之一。其推论之一是：余伴随表示带有所谓的 Kostant--Kirillov Poisson 括号，所以我们预期这里会出现某种 Poisson 几何。一种说法是：
\begin{equation}\tag{16.16}
\cM^1_B=T^*[3](\mathrm{pt}/G^\vee).
\end{equation}
这是 $3$-移位辛的，即一个 $\mathbb{P}_4/\mathbb{E}_4$ 叠。要点在于切空间是
\begin{equation}\tag{16.17}
T(\mathrm{pt}/G^\vee)=\fg^\vee[1],
\end{equation}
于是
\begin{equation}\tag{16.18}
T^*(\mathrm{pt}/G^\vee)=\fg^\vee[-1].
\end{equation}
它被 $3$-移位，所以定理中那个 $2$-移位空间实际上是按 $-1+3$ 移位的。% [Tutor Note: 这里作者把两种描述混在一起：$\cM^1_{B_{G^\vee}}=\fg^{\vee*}[2]/G^\vee$ 按定义是 $2$-移位辛的；而 $T^*[3](\mathrm{pt}/G^\vee)$ 是 $3$-移位辛的。两者的关系是 $\fg^{\vee*}[2]/G^\vee\simeq T^*[3](\mathrm{pt}/G^\vee)$：切空间 $\fg^\vee[1]$ 的余切移 $[2]$ 得 $\fg^{\vee*}[2]$，而 $\mathrm{pt}/G^\vee$ 的余切丛本身已带 $3$-移位（因为 $T(\mathrm{pt}/G^\vee)$ 在次数 $1$，余切在 $-1$，再加 $3$）。作者所说“$-1+3$”即指这个抵消。]

\subsection{它从何而来？}

考虑 $\Loc_{G^\vee}(S^2)$。我们知道 $S^2$ 可以描述为两个圆盘沿圆周粘合。由于圆盘可缩，圆盘上的局部系统实际上是平凡的：圆盘上只有一个局部系统，其自同构由 $G^\vee$ 给出。即
\begin{equation}\tag{16.19}
\Loc_{G^\vee}(D^2)=\mathrm{pt}/G^\vee.
\end{equation}
于是我们有两份 $\mathrm{pt}/G^\vee$，并想沿圆周上的局部系统粘合。但圆周上确有有意思的局部系统，因为可以有单值性（monodromy）。事实上，圆周上的局部系统由它的单值性在共轭意义下完全决定：
\begin{equation}\tag{16.20}
\Loc_{G^\vee}(S^1)=G^\vee/G^\vee.
\end{equation}
把这些放在一起得到：
\begin{equation}\tag{16.21}
\Loc_{G^\vee}(S^2)=(\mathrm{pt}/G^\vee)\times_{G^\vee/G^\vee}(\mathrm{pt}/G^\vee).
\end{equation}

\begin{srcrem}{115}
我们隐含地用了：余极限上的局部系统可以描述为各块上局部系统的极限。这是同义反复，因为可以把 $\Loc$ 看成
\begin{equation}\tag{16.22}
\Loc_{G^\vee}(-)=\Hom(-,G^\vee).
\end{equation}
\end{srcrem}

我们可以把它改写成\footnote{确实需要想一下为什么商与交可以交换。}
\begin{align}
\Loc_{G^\vee}(S^2)&=(\mathrm{pt}/G^\vee)\times_{G^\vee/G^\vee}(\mathrm{pt}/G^\vee)\tag{16.23}\\
&=(\mathrm{pt}\times_{G^\vee}\mathrm{pt})/G^\vee.\tag{16.24}
\end{align}

也就是说，$S^2$ 上的局部系统可以描述为 $G^\vee$ 中一点的自交，再模去 $G^\vee$。在阿贝尔情形，一点的自交是平凡的；但在导出情形不是。这就是所谓的\term{基于一点的环路空间}{based loop space}\footnote{要注意记号与语境：这是导出代数几何意义下的基于一点的环路，而不是同伦论中的环路空间。}，记作
\begin{equation}\tag{16.25}
\Omega G^\vee=\mathrm{pt}\times_{G^\vee}\mathrm{pt},
\end{equation}
即
\begin{equation}\tag{16.26}
\Loc_{G^\vee}(S^2)=(\mathrm{pt}\times_{G^\vee}\mathrm{pt})/G^\vee=\Omega G^\vee/G^\vee.
\end{equation}
这些环路“非常小”：它们探测不到 $G^\vee$ 的整体拓扑。事实上，它们只依赖于单位元的一个小邻域。所以若选取一些坐标，可以等同
\begin{equation}\tag{16.27}
\Loc_{G^\vee}(S^2)=(\mathrm{pt}\times_{\fg}\mathrm{pt})/G^\vee.
\end{equation}
也就是说，我们是在计算向量空间中一点的导出自交。% [Tutor Note: 这里的 $\mathrm{pt}\times_{\fg}\mathrm{pt}$ 中，$\fg$ 应理解为 $G^\vee$ 在单位元处的形式邻域（或直接取 $\fg$ 作为加群）；(16.27) 的等同只在单位元的无穷小邻域上成立，这正是“环路非常小”的含义。与 (16.25) 的差别在于把群换成了它的李代数。]

导出世界的全部要点，就是修正纤维积（特别是交）在非横截时的行为。设 $V$ 是向量空间，考虑 $0\in V$ 的导出自交：
\begin{equation}\tag{16.28}
0\cap_V0=\mathrm{pt}\times_V\mathrm{pt},
\end{equation}
即 $V$ 在 $0$ 处的基于一点的环路空间。这是导出交，但我们从记号中省略了“导出”二字。这个自交在横截性方面是“最坏的情形”：$0$ 的切空间是平凡的，所以当我们问这个交是否横截（即切空间是否张成环境空间在该点的切空间）时，会发现我们缺失了整个切空间 $T_0V\simeq V$。所以这离横截要多远有多远，但导出结构恰好会记录这一退化。

照代数几何的惯例，我们将通过其上的函数来理解这个空间。我们用到的原则是：取 $\cO$ 时，交变成张量积：
\begin{equation}\tag{16.29}
\cO(\mathrm{pt}\times_V\mathrm{pt})=k\otimes_{\cO(V)}k=k\otimes_{\mathrm{Sym}V^*}k,
\end{equation}
因为 $\cO(V)=\mathrm{Sym}V^*$。这是导出张量积，但省略记号。用 Koszul 消解可以算出这是外代数：
\begin{equation}\tag{16.30}
\cO(\mathrm{pt}\times_V\mathrm{pt})=k\otimes_{\cO(V)}k\simeq\wedge^\bullet V^*.
\end{equation}
这是一个交换（微分）分次代数（commutative (d)ga）：
\begin{equation}\tag{16.31}
\wedge^\bullet V^*=\mathrm{Sym}(V^*[1]).
\end{equation}
现在我们要计算它上的模范畴，即这个导出交 $\mathrm{pt}\times_V\mathrm{pt}$ 上的拟凝聚层。无论这个模范畴是什么，先指出它有一些结构。它会有张量积，因为它是交换环上的模；但我们对这个乘法不感兴趣。我们感兴趣的是（结合的）卷积结构。

每当我们有来自映射 $X\to Y$ 的 $X\times_YX$ 时，它就是某种以两种方式映入 $X$ 的群胚。它有卷积结构，例如当 $X$ 是有限集时，这个卷积结构就是矩阵乘法。在我们这里，$X=\mathrm{pt}$，$Y=V$。无论这个卷积结构是什么，它在 Koszul 对偶另一侧对应的运算会更明显。

为描述范畴 $\wedge^\bullet V^*\text{-}\mathrm{mod}$，我们将聚焦于我们真正理解的仅有的对象：自由模与增广模。增广模就是 $k$，其中 $\wedge^\bullet V^*\text{-}\mathrm{mod}$ 的所有外代数生成元都作用为零。它是卷积的单位。

另一种观点是：这个导出交是 $0$ 处的基于一点的环路空间：
\begin{equation}\tag{16.32}
\mathrm{pt}\times_V\mathrm{pt}=\Omega_0V.
\end{equation}
我们知道基于一点的环路空间在环路拼接下构成群（至少在某种同伦意义下）。其单位是单位元处的摩天层，即增广模。

为理解这个范畴，我们的策略如下。要研究范畴 $\cC$，可以考虑一个我们喜欢的对象 $P$，然后希望该范畴可以被描述为 $A=\Hom(P,P)$ 上的模。也就是说，有一个函子
\begin{equation}\tag{16.33}
\cC\xrightarrow{\ \Hom(P,-)\ }(\Hom(P,P))\text{-}\mathrm{mod},
\end{equation}
在好的情形下它是等价。

在我们的情形中，我们将通过计算 $\End(k)$ 来理解 $\wedge^\bullet V^*\text{-}\mathrm{mod}$。这将把我们引向 Koszul 对偶。所以需要计算：
\begin{equation}\tag{16.34}
\End_{\wedge^\bullet V^*}(k).
\end{equation}
有趣的是，$\mathrm{pt}\times_V\mathrm{pt}$ 是奇异的，即它是某种奇怪且导出的东西。如我们之前看到的，若取奇点处的摩天层，则它有无穷维自同态。也就是说，我们可以通过注意到 $k\in\wedge^\bullet V^*\text{-}\mathrm{mod}$ 不是完美的，来察觉自己在奇点处。对此的解释是：这看起来像奇异支撑上的函数，即空间在余切丛中奇点处粘上的那些额外“翅膀”上的函数。回想图 15.1 与 15.2 中的例子：节点在奇点处长出一只翅膀，用来解释摩天层的巨大自同态。这不是同一个例子，但现象相同。这就是为什么我们应当期望出现类似对称代数的东西：
\begin{equation}\tag{16.35}
\End_{\wedge^\bullet}(k)\simeq S=\mathrm{Sym}(V[-2])=\cO(V^*[2]).
\end{equation}
于是我们的空间长出了一只“翅膀”：$V^*[2]$。% [Tutor Note: 原文此处页边有手写批注“Good reference: Goresky--Kottwitz--Macpherson”。这批注指向 GKM 关于等变上同调与环面作用的经典工作；在下文的类比中，$S=\mathrm{Sym}(V[-2])$ 正是 $S^1$ 等变上同调 $H^*_{S^1}(\mathrm{pt})=k[u]$（$|u|=2$），见 (16.38)。]

这里代数内容很多。我们不会在拓扑情形下解释它。我们将把这个外代数看成：
\begin{equation}\tag{16.36}
\wedge^\bullet V^*=H_*(S^1).
\end{equation}
回想 $H_*(S^1)$ 是一个生成元上的外代数。若 $V$ 不是一维的，就把 $S^1$ 换成一个环面来想象。$(H_*(S^1),*)$ 是一个结合代数，可以称为同伦论中 $S^1$ 的群代数。对有限群，群同调字面上就是群代数。一般地，群的同调总带有结合乘法，称为 Pontrjagin 乘积。于是我们有
\begin{equation}\tag{16.37}
S=\End_{\wedge^\bullet}(k)=\End_{H_*(S^1)\text{-}\mathrm{mod}}(H^*(\mathrm{pt}))=\Hom_{S^1}(k,H^*(\mathrm{pt})).
\end{equation}
这是等变上同调的另一种说法：
\begin{equation}\tag{16.38}
S=H_{S^1}(\mathrm{pt})=H^*(BS^1\simeq\C\PP^\infty)\simeq k[u],
\end{equation}
其中 $|u|=2$。

Koszul 对偶大致是说：
\begin{equation}\tag{16.39}
\wedge^\bullet V^*\text{-}\mathrm{mod}\xrightarrow[\ \sim\ ]{\Hom_{\wedge^\bullet V^*}(k,-)}S\text{-}\mathrm{mod}.
\end{equation}
这里会遇到一些标准的技术困难。首先，必须小心我们考虑的是哪类模。例如我们可能考虑凝聚的 $\wedge^\bullet V^*$-模 $\mathrm{Coh}_\Lambda$（即有有限维上同调的有界复形），于是有范畴等价
\begin{equation}\tag{16.40}
\mathrm{Coh}_\Lambda\leftrightarrow\mathrm{Coh}_S.
\end{equation}
有趣的是，$S$ 只是对称代数，所以它看起来像向量空间上的层，而向量空间是光滑的，于是
\begin{equation}\tag{16.41}
\mathrm{Coh}_S\simeq\mathrm{Perf}_S.
\end{equation}
这非常漂亮，只不过
\begin{equation}\tag{16.42}
\mathrm{Coh}_\Lambda\neq\mathrm{Perf}_\Lambda.
\end{equation}
问题在于，增广模 $k\in\mathrm{Coh}_\Lambda$ 被送到
\begin{equation}\tag{16.43}
k\mapsto\Hom_{\wedge^\bullet V^*}(k,k)=S.
\end{equation}
但 $k$ 不是完美/紧的，而 $S\in\mathrm{Coh}_S$ 是。这一现象的另一个版本如下。我们有 $k[u,u^{-1}]\in\mathrm{Perf}_S$，它在（伴随函子下）被送到 $0\in\mathrm{Perf}_\Lambda$。所以这里有某种张力，而它由ind-凝聚层（ind-coherent sheaves）来解决。记作 $\mathrm{IndCoh}$ 或 $\QC^!$。其想法是：看奇异空间时，我们加一个叹号来提醒你得小心(!)。但要点是：
\begin{equation}\tag{16.44}
\QC^!(\mathrm{pt}\times_V\mathrm{pt})\simeq S\text{-}\mathrm{mod}.
\end{equation}
我们应把它看成对 $\Lambda\text{-}\mathrm{mod}$ 的修正/扩大：
\begin{equation}\tag{16.45}
\QC^!(\mathrm{pt}\times_V\mathrm{pt})\supset\Lambda\text{-}\mathrm{mod}=\QC(\mathrm{pt}\times_V\mathrm{pt}).
\end{equation}
其想法是：$k$ 在凝聚世界里不是紧的（行为不良），但当我们过渡到导出凝聚世界时，$k$ 变成一个紧对象（行为良好）。

无论如何，Koszul 对偶是如下形式的陈述：
\begin{equation}\tag{16.46}
(\Lambda\text{-}\mathrm{mod},*)\xrightarrow{\ \sim\ }(S\text{-}\mathrm{mod},\otimes).
\end{equation}
我们可以把它看作 Cartier/Pontrjagin 对偶的一个版本。这告诉我们：$\wedge^\bullet V^*\text{-}\mathrm{mod}$ 上的卷积结构，在通过 Koszul 对偶之后对应于一个自然出现的张量积结构。

回想在导出 Satake 的语境中，我们感兴趣的空间是
\begin{equation}\tag{16.47}
\Loc_{G^\vee}(S^2)=(\mathrm{pt}\times_{\fg^\vee}\mathrm{pt})/G^\vee.
\end{equation}

我们感兴趣的范畴是
\begin{equation}\tag{16.48}
\mathcal{B}_{G^\vee}(S^2)=\QC^!(\Loc_{G^\vee}(S^2))=\QC^!(\mathrm{pt}\times_{\fg^\vee}\mathrm{pt})^{G^\vee}.
\end{equation}
我们用 $\QC^!$ 来提醒自己正在处理奇异空间，因而把奇点处的摩天层当作“受尊重的社会成员”。回想阿贝尔版本考虑的是 $\QC(\mathrm{pt})^{G^\vee}$（或 $\QC^!$，但因为 $\mathrm{pt}$ 并不奇异，二者没有区别）；但在导出版本中我们看到，实际上应当考虑点上的导出自交上的层，而不是点本身。回到计算，我们有
\begin{equation}\tag{16.49}
\mathcal{B}_{G^\vee}(S^2)=(\mathrm{Sym}(\fg^\vee[-2]))^{G^\vee},
\end{equation}
并带张量积。对偶化后我们看到，这是（$G^\vee$-等变的）层，位于
\begin{equation}\tag{16.50}
(\fg^{\vee*}[2])/G^\vee=T^*[3](\mathrm{pt}/G^\vee).
\end{equation}

\section{局部算子}
\begin{sectionkey}
本节把“线算子的导出版本”落到局部算子上：$A$-侧局部算子为 $H^*_G(\mathrm{pt})=\mathrm{Sym}\,\mathfrak{h}^*[2]^W$，$B$-侧为 $(\mathrm{Sym}(\fg^\vee[-2]))^{G^\vee}=\mathrm{Sym}(\mathfrak{h}^\vee[-2])^W=\cO(\mathfrak{h}^\vee[2]/W)$。两者通过特征多项式映射 $\cM^1=\fg^{\vee*}[2]/G^\vee\to\cM^0=\mathfrak{h}^{\vee*}[2]/W$ 相联系，$\cM^0$ 即 Coulomb 分支。$B$-侧的关键是：这些局部算子通过“测量奇异支撑的特征多项式”作用在 $\QC^!(\Loc_{G^\vee})$ 上；限制到该映射的 $0$-纤维，就得到 Arinkin--Gaitsgory 的幂零奇异支撑范畴。
\end{sectionkey}

回想在阿贝尔故事中我们只有平凡的局部算子。让我们重新审视这一点。回想局部算子由单位线算子的自同态给出。所以这里它是 $T^*[3](\mathrm{pt}/G^\vee)$ 的仿射化。也就是说，它是
\begin{equation}\tag{16.51}
(\mathrm{Sym}(\fg^\vee[-2]))^{G^\vee}=\mathrm{Sym}(\mathfrak{h}^\vee[-2])^W=\cO(\mathfrak{h}^\vee[2]/W),
\end{equation}
其中 $\mathfrak{h}$ 是 Cartan 子代数。所以它是李代数上通常的不变量多项式空间的分次版本。换言之，我们有
\begin{equation}\tag{16.52}
\cM^1=\fg^{\vee*}[2]/G^\vee,
\end{equation}
然后还有 $\cM^0$，它通常称为（$\mathcal{N}=4$ SYM 的）\term{Coulomb 分支}{Coulomb branch}：
\begin{equation}\tag{16.53}
\cM^0=\mathfrak{h}^{\vee*}[2]/W,
\end{equation}
它们之间的映射由特征多项式映射给出：
\begin{equation}\tag{16.54}
\cM^1=\fg^{\vee*}[2]/G^\vee\to\cM^0=\mathfrak{h}^{\vee*}[2]/W.
\end{equation}
所以我们确实找到了局部算子：它们看起来正是特征多项式的系数。它们全都位于非零次数，这就是我们之前没看到它们的原因。% [Tutor Note: 这里用的是 Chevalley 限制定理：$(\mathrm{Sym}\,\fg^\vee)^{G^\vee}\cong(\mathrm{Sym}\,\mathfrak{h}^\vee)^W$，即不变多项式由其在 Cartan 上的限制决定；特征多项式映射 $\fg^{\vee*}\to\mathfrak{h}^{\vee*}/W$ 正是这一限制的几何实现。]

\subsection{$A$-侧}

既然我们有了一些非平凡的局部算子，就可以问一些问题：在实践中去哪里找它们，它们探测/做什么？这在自守侧最容易看出来。考虑 Cartan 模去 Weyl 群：
\begin{equation}\tag{16.55}
\mathfrak{h}^{\vee*}[2]/W=\mathfrak{h}[2]/W.
\end{equation}
通常方便地就把它写成 $\mathfrak{h}/W$，并只用 $\C^\times$ 作用来追踪分次。这个空间以非常自然的方式出现。回想自守 Satake 范畴是
\begin{equation}\tag{16.56}
\mathcal{A}_G(S^2)=\Shv(LG_+\backslash LG/LG_+),
\end{equation}
即仿射 Grassmann 簇上 $LG_+$-等变的层。这是某种导出怪兽，但我们只需关心卷积下的单位，即恒等陪集处的摩天层：我们有 $\mathrm{pt}\in\Gr=LG/LG_+$，并且看的是 $LG_+$-等变的层，所以单位是摩天层 $k_{\mathrm{pt}/LG_+}$。但同伦地，这个 $LG_+$ 作用与 $G$-作用相同：
\begin{equation}\tag{16.57}
k_{\mathrm{pt}/LG_+}\simeq k_{\mathrm{pt}/G}.
\end{equation}
特别地，这意味着 $A$-侧的局部算子构成
\begin{equation}\tag{16.58}
\mathcal{A}_G(S^3)=\End(1)=H^*(BG)=H^*_G(\mathrm{pt}).
\end{equation}
如同在拓扑中，
\begin{equation}\tag{16.59}
H^*_G(\mathrm{pt})=\mathrm{Sym}\,\mathfrak{h}^*[2]^W\simeq(\mathrm{Sym}\,\fg^*[2])^G.
\end{equation}
这是用冗长的方式说：$A$-侧的局部算子就是 $H^*_G(\mathrm{pt})$，即 $G$-等变上同调的系数环。这是诸如 Donaldson 理论中熟悉的东西：自然的系数环由等变上同调给出。

回想局部算子的要点是：在 $A$-侧一切都线性于它。也就是说，像 $\Shv(\Bun_G(C))$ 这样的东西可以写成对 $G$ 的商：
\begin{equation}\tag{16.60}
\Shv(\Bun_G(C))=\Shv(\Bun_G(C,x)/G),
\end{equation}
意思是它线性于 $H^*(BG)$。

\subsection{$B$-侧}

$B$-侧局部算子与线算子的导出加厚将测量奇异支撑。这正是 Arinkin--Gaitsgory 所描述的。设 $C$ 是 Riemann 面（或者干脆只是一个拓扑曲面）。$C$ 上的状态构成
\begin{equation}\tag{16.61}
\mathcal{B}_{G^\vee}(C)=\QC^!(\Loc_{G^\vee}(C)).
\end{equation}
这类层的范畴有不同模型。例如我们过去总写 $\QC$（拟凝聚层），但现在我们意识到必须更小心，要用某种ind-凝聚层，即 $\QC^!$，因为这是奇异空间。如上次所见，奇异空间上的层有奇异支撑的概念。特别地，这个空间 $\Loc_{G^\vee}(C)$ 作为纤维积出现：
\begin{equation}\tag{16.62}
\begin{array}{ccc}
\Loc_{G^\vee} & \longrightarrow & (G^\vee)^{2g}/G^\vee\\[2pt]
\downarrow & & \downarrow{\scriptstyle\Pi_{[,]}}\\[-2pt]
1/G^\vee & \longrightarrow & G^\vee/G^\vee
\end{array}
\end{equation}
它将带有来自基于一点的环路的对称性：
\begin{equation}\tag{16.63}
\Omega_1G^\vee/G^\vee\ \circlearrowleft\ \Loc_{G^\vee}.
\end{equation}
要点是：尽管 $\Loc_{G^\vee}(C)$ 不光滑，它却非常接近光滑：它是两个光滑空间之间映射的纤维。这个“翅膀”的空间，即奇异支撑所居住的地方 $T^*[1]\Loc_{G^\vee}$，看起来其实是：一个局部系统 $E$，连同 $\fg^{\vee*}[2]$ 的一个被 $E$ 扭曲的截面。无论这东西是什么，我们都可以只关心不变多项式部分：
\begin{equation}\tag{16.64}
T^*[1]\Loc_{G^\vee}\to\mathfrak{h}^{\vee*}[2]/W.
\end{equation}
这一切只是说：我们所说的局部算子 $\cO(\mathfrak{h}^{\vee*}[2])^W$ 通过测量（特征多项式意义下的）奇异支撑，作用在范畴 $\QC^!(\Loc)$ 上。

Arinkin--Gaitsgory 提出限制到这个映射的 $0$-纤维，即具有\term{幂零奇异支撑}{nilpotent singular support}的ind-凝聚层，记作
\begin{equation}\tag{16.65}
\QC^!_N(\Loc_{G^\vee}).
\end{equation}
他们这样提议，是因为它与 $A$-侧相匹配。$\Shv(\Bun_G)$ 的“显然”版本是 $H^*_G(\mathrm{pt})$-挠的。也就是说，它生活在
\begin{equation}\tag{16.66}
0\in\mathfrak{h}[2]/W=\cM^0.
\end{equation}
所以要把它与谱侧的东西匹配起来：
\begin{equation}\tag{16.67}
\Shv(\Bun_G)\stackrel{?}{\simeq}\QC^!_N(\Loc_{G^\vee}),
\end{equation}
它也应该生活在 $0$ 处。

\begin{srcthm}{17}[Beraldo]
由线算子构造的 Riemann 面 $C$ 上的可观测量构成
\[
\int_C\mathcal{B}_{G^\vee}(S^2)\simeq\mathrm{DQ}(T^*[1]\Loc_{G^\vee}),
\]
其中 $\mathrm{DQ}$ 表示“形变量子化模”，即这个空间量子化上的层。
\end{srcthm}

这是说：若用线算子作为可观测量，就能完全探测凝聚层的奇异支撑：线算子不多不少正好探测奇异支撑。可以与之前的比较：
\begin{equation}\tag{16.68}
\int_C\Rep(G^\vee)=\QC(\Loc_{G^\vee}(C)).
\end{equation}
这些是 $C$ 上“阿贝尔”线算子给出的可观测量，所以它自然地作用在下面这个范畴上：
\begin{equation}\tag{16.69}
\mathcal{A}_G(C)=\Shv(\Bun_G(C)).
\end{equation}
于是当我们从线算子过渡到导出线算子时，得到了奇异支撑。

早该说的是：导出几何 Satake 断言
\begin{equation}\tag{16.70}
\Shv(LG_+\backslash LG/LG_+)\simeq\QC^!(\Loc_{G^\vee}(S^2))=(\mathrm{Sym}(\fg^\vee[-2])\text{-}\mathrm{mod})^{G^\vee}.
\end{equation}
这是 Bezrukavnikov--Finkelberg 的定理。也就是说，导出几何 Satake 说导出线算子在 $A$-侧与 $B$-侧相同。% [Tutor Note: 原文此处有手写批注“cite”，提示补引 Bezrukavnikov--Finkelberg 关于导出几何 Satake 的工作。该定理通常表述为 $\Shv(LG_+\backslash LG/LG_+)\simeq\QC^!(\Loc_{G^\vee}(S^2))$，等价于 $\fg^{\vee*}[2]/G^\vee$ 上的ind-凝聚层；右端的 $(\mathrm{Sym}(\fg^\vee[-2])\text{-}\mathrm{mod})^{G^\vee}$ 是对偶描述。]

所以要点是：由于 $\Loc_{G^\vee}C$ 奇异，线算子的导出版本在谱意义上分解于 $T^*[1]\Loc_{G^\vee}$ 之上。
% [Tutor Note: 原文 "the derived version of line operators spectrally composes over $T^*[1]\Loc_{G^\vee}$
%  里的 composes 不好直译。作者在全书表述这类谱分解时一律用 spectrally decomposes
%  （第 1、3、10、11、13、14、15 章共 14 处），此处 composes 应是 decompose 的笔误或口语变体。
%  译文按 "分解于 $T^*[1]\Loc_{G^\vee}$ 之上" 处理，即导出线算子按该处的奇异支撑来组织。
%  [译者按（无外部依据）：依据是作者他处的一致用法与本章语境，未找到外部文献佐证。]]

\begin{srcrem}{116}
光滑点就是所谓的\term{尖点 Langlands 参数}{cuspidal Langlands parameters}。所以整套奇异支撑理论告诉你如何理解非尖点的 Langlands 参数。
\end{srcrem}

\section{边界条件（余维一缺陷）}
\begin{sectionkey}
本节讨论余维一缺陷：它们把流形分成两部分，因而自然地充当两个场论之间的“映射”。把两张理论拼起来的畴壁，可以通过沿畴壁“折叠”而看成单一理论 $Z\otimes Z'$ 的边界条件。由此，Langlands 函子性（群同态 $H^\vee\to G^\vee$ 诱导 $B_{H^\vee}\to B_{G^\vee}$ 的畴壁）成为边界条件的一个特例。最后的结论是：$B_{G^\vee}$ 的边界条件大致对应于（移位的）哈密顿 $G^\vee$-空间。
\end{sectionkey}

到目前为止我们谈了点（局部）算子、线算子和面算子。现在来谈余维一算子/缺陷。我们可以考虑两类不同的余维一缺陷。与早先的缺陷概念不同，余维一缺陷自然把流形分成两部分。你可以谈边界，或嵌入的余维一子流形。在物理中它们被称为\term{畴壁}{domain walls}。因为空间被分成两部分，畴壁一般就是把两个不同的场论 $\cZ$ 与 $\cZ'$ 分隔开的东西。场论之间最自然的映射概念之一就是畴壁，即应把分隔两个理论的畴壁看成从一个理论到另一个理论的映射。形式地说，我们把定义在通常配边上的理论推广到定义在二分（bipartite）流形上。也就是说，从 $\cZ$ 到 $\cZ'$ 的映射，是把两者都延拓到二分流形上定义的，其中一侧是 $\cZ$，另一侧是 $\cZ'$。

\subsection{Langlands 函子性}

我们不打算把 $A$-侧与 $B$-侧的理论联系起来，尽管可以这样想：Langlands 对应本身就由某个神奇的畴壁给出。相反，我们将做下面这件事。谈 Langlands 函子性时，人们说的是：对两个不同的群 $G^\vee$ 与 $H^\vee$，我们可能试图把 $\mathcal{B}_{G^\vee}$ 与 $\mathcal{B}_{H^\vee}$ 联系起来。最显然的做法（这也正是人们通常所说的 Langlands 函子性）是取一个子群 $H^\vee\subset G^\vee$，然后构造映射
\begin{equation}\tag{16.71}
\mathcal{B}_{H^\vee}\to\mathcal{B}_{G^\vee}.
\end{equation}
特别地，若看 Riemann 面 $C$ 上的局部系统，就得到一个真正的映射
\begin{equation}\tag{16.72}
\Loc_{H^\vee}(C)\to\Loc_{G^\vee}(C),
\end{equation}
它诱导出前推函子
\begin{equation}\tag{16.73}
\QC^!(\Loc_{H^\vee}(C))\to\QC^!(\Loc_{G^\vee}(C)).
\end{equation}
我们没有用到 $C$ 是 Riemann 面。事实上这可以延拓到整个场论。我们甚至没有用到它是个子群，所以完全可以从群同态 $H^\vee\to G^\vee$ 出发。于是 Langlands 函子性问的是：这在 $A$-侧对应什么：
\begin{equation}\tag{16.74}
\begin{array}{ccc}
\mathcal{B}_{H^\vee} & \longrightarrow & \mathcal{B}_{G^\vee}\\[2pt]
\wr & & \wr\\[2pt]
\mathcal{A}_H & \stackrel{?}{\dashrightarrow} & \mathcal{A}_G
\end{array}
\end{equation}
% [Tutor Note: (16.74) 是 Langlands 函子性的图示：中间一行的两个 $\wr$ 表示上下两行对象在几何 Langlands 对偶下相互对应（左侧 $H$-那一对、右侧 $G$-那一对），底行 $\mathcal{A}_H$ 到 $\mathcal{A}_G$ 的虚线加“?”表示这个函子既未构造、也未必存在——正如紧接的正文所说，$H^\vee\to G^\vee$ 并不给出任何形式的映射 $H\to G$。[译者按（无外部依据）：依 (16.71)–(16.73) 与下一句推得。]]
问题在于，$H^\vee\to G^\vee$ 并不给出任何形式的映射 $H\to G$。

这实际上更接近 Langlands 最初表述纲领的方式。也就是说，存在某个 Langlands 函子性原则，它不精确地说：对任何映射 $H^\vee\to G^\vee$，都应存在一个从 $H$ 的自守形式到 $G$ 的自守形式的操作。% [Tutor Note: 原文明明讨论的是 automorphic forms（自守形式），但抽取文本两处都写作 “automorphism forms”（自同构形式），这是笔误；译文按“自守形式”处理。]

这类函子性的另一个关键例子是抛物诱导/Eisenstein 级数理论。在这个语境中，给定群 $G^\vee$，然后取极大环面 $T^\vee$。

\begin{figure}[htbp]
  \centering
  \includegraphics[width=0.72\textwidth]{figure16-1.png}
  \caption{沿畴壁“折叠”后，我们可以把 $\cZ$ 与 $\cZ'$ 之间的畴壁看作张量积 $\cZ\otimes\cZ'$。（原文 p.~275，Figure 16.1）}
  \label{fig:16-1}
\end{figure}

我们可以把它看成一个子群，但这里我们改为把 Borel 子群看成在这两者之间进行插值：
\begin{equation}\tag{16.75}
\begin{array}{ccc}
& B^\vee &\\
\swarrow & & \searrow\\
G^\vee & & T^\vee
\end{array}
\end{equation}
这给出 $\mathcal{B}_{G^\vee}$ 与 $\mathcal{B}_{T^\vee}$ 之间的畴壁。例如在 Riemann 面 $C$ 上，我们得到一个对应：
\begin{equation}\tag{16.76}
\begin{array}{ccc}
& \Loc_{B^\vee}(C) &\\
\swarrow & & \searrow\\
\Loc_{G^\vee}(C) & & \Loc_{T^\vee}(C)
\end{array}
\end{equation}
$\Loc_{B^\vee}(C)$ 由带旗的 $G^\vee$-局部系统构成。于是有一个自然的操作，从 $\Loc_{G^\vee}$ 得到 $\Loc_{T^\vee}$：拉回与前推。在这种情形下，我们其实知道另一侧该怎么做：我们同样想
\begin{equation}\tag{16.77}
G\supset B\twoheadrightarrow T,
\end{equation}
于是得到对应
\begin{equation}\tag{16.78}
\begin{array}{ccc}
& \Bun_B &\\
\swarrow & & \searrow\\
\Bun_G & & \Bun_T
\end{array}
\end{equation}
所以我们又可以拉回与前推。

\subsection{畴壁与边界条件}

现在稍微换一个视角。双模 ${}_AM_B$ 等同于 $A\otimes B^{\mathrm{op}}$ 上的模。在线性代数中，这只是说映射 $f:V\to W$ 等同于元素 $f\in V^*\otimes W$。所以我们取双模 ${}_AM_B$，可以把它看成一个函子
\begin{equation}\tag{16.79}
{}_AM_B:A\text{-}\mathrm{mod}\to B\text{-}\mathrm{mod},
\end{equation}
并等价地把它看成 $A\otimes B^{\mathrm{op}}\text{-}\mathrm{mod}$ 的元素。

直观地说，与其考虑双模/畴壁，我们将像图 16.1 那样把纸张折叠起来。边界条件就是到平凡理论的畴壁。所以，与其把余维一缺陷看成从 $\cZ$ 到 $\cZ'$ 的畴壁/映射，我们把它看成单一理论 $\cZ\otimes\cZ'$ 的边界条件。要点是：畴壁（或函子性）是边界条件的特例。

例如，在这个 Eisenstein 级数的例子里，与其把畴壁看成由 $B$ 给出：
\begin{equation}\tag{16.80}
G\supset B\twoheadrightarrow T,
\end{equation}
我们可以把它折过来，把 $B$ 看成 $G\times T$ 理论的边界条件。或者可以几何地看。若有
\begin{equation}\tag{16.81}
U\hookrightarrow B\twoheadrightarrow T,
\end{equation}
则 $G/U$ 是一个带有两个作用的空间：
\begin{equation}\tag{16.82}
G\ \circlearrowleft\ G/U\ \circlearrowright\ T.
\end{equation}
然后对群取商，就得到
\begin{equation}\tag{16.83}
\begin{array}{ccc}
& \mathrm{pt}/B &\\
\swarrow & & \searrow\\
\mathrm{pt}/G & & \mathrm{pt}/T
\end{array}
\end{equation}
要点是：我们可以只关注边界条件，而不必操心畴壁。

\subsection{哈密顿作用}

要点将是下面这个。

\begin{srcclaim}{3}
$\mathcal{B}_{G^\vee}$ 的边界条件大致等同于/给出一个（移位的）哈密顿 $G^\vee$-空间。
\end{srcclaim}

这是 Coulomb 分支理论（Braverman--Finkelberg--Nakajima）的一个版本，它早先已在物理语境中出现，例如 Gaiotto--Witten。% [Tutor Note: 原文此处页边有手写批注“cite BFN and GW”，提示补引 Braverman--Finkelberg--Nakajima（Coulomb 分支）与 Gaiotto--Witten（边界条件与规范理论）的文献。译文正文保留了对这两项工作的提及。]

\begin{srcex}{112}
设 $G$ 是群，$X$ 是 $G$-空间。现在考虑 $T^*X$。它是辛的，并从 $X$ 得到 $G$ 的作用，但它还有一个矩映射
\begin{equation}\tag{16.84}
T^*X\to\fg^*,
\end{equation}
它是 $G$-等变的。这个矩映射给出 $\mathrm{Lie}(G)=\fg$-作用的哈密顿量。其想法是，映射 $T^*X\to\fg^*$ 给出函数上的映射
\begin{equation}\tag{16.85}
H:\cO(\fg^*)=\mathrm{Sym}\,\fg\to\cO(T^*X),
\end{equation}
即哈密顿算子。于是对 $x\in\fg$，可以看对应的哈密顿函数 $H_x\in\cO(T^*X)$；说它对作用而言是哈密顿的，就是说与 $H_x$ 的 Poisson 括号等同于 $x$ 的作用：
\begin{equation}\tag{16.86}
\{H_x,-\}=\text{action of }x\text{ on }T^*X.
\end{equation}
\end{srcex}

于是我们有一个映射 $\fg\to\cO(T^*M)$，它提升 $\fg$ 到向量场的映射：
\begin{equation}\tag{16.87}
\begin{array}{ccc}
\mathrm{Sym}\,\fg & \xrightarrow{\ H\ } & \cO(T^*M)\\[2pt]
\uparrow{\scriptstyle\iota} & & \downarrow{\scriptstyle\{,-\}}\\[2pt]
\fg & \xrightarrow{\ \text{action}\ } & \mathrm{Vect}(T^*M)
\end{array}
\end{equation}
% [Tutor Note: 原图是一幅不太规整的草图：$\mathrm{Sym}\,\fg$ 有斜箭头分别指向 $\cO(T^*M)$ 与 $\mathrm{Vect}(T^*M)$，$\fg$ 有箭头指向 $\mathrm{Vect}(T^*M)$。这里按数学含义重排为交换方块：上横为 (16.85) 的 $H$，右竖为“取哈密顿向量场”（Poisson 括号），下横为 $G$-作用，左竖为 $\fg\hookrightarrow\mathrm{Sym}\,\fg$。]

现在若给定任意带 $G$-作用的辛空间 $M$（即不一定形如 $T^*X$），我们可以要求同样的数据：
\begin{equation}\tag{16.88}
\begin{array}{ccc}
M & \longrightarrow & \fg^*\\[2pt]
\circlearrowleft & {\scriptstyle G} & \circlearrowright
\end{array}
\end{equation}
即一个到 $\fg^*$ 的映射，对群等变，生成 $G$ 的李代数作用，且与辛结构相容。

我们已经看到某种应当让我们想起几何 Satake 的东西。我们可以把这组数据重新表述为叠的映射：
\begin{equation}\tag{16.89}
M/G\xrightarrow{\ \mu\ }\fg^*/G.
\end{equation}
除去分次问题（例如我们或许会期望 $\fg^*[2]/G$），矩映射的目标看起来像 $B$-理论的真空模空间：
\begin{equation}\tag{16.90}
M/G\xrightarrow{\ \mu\ }\fg^*[2]/G=\cM^1_B.
\end{equation}
所以哈密顿几何切入的地方在于：到 $\cM^1_B$ 的映射会非常接近哈密顿 $G$-作用。因此我们希望从边界条件中提取出看起来像哈密顿作用的数据。

让我们再做些场论缺陷学。想法是做同样的谱废话。我们这样总结：给定 TFT $\cZ$ 的一个边界条件 $T$，我们想提取算子代数并取 $\Spec$。这将得到空间 $M$，而在我们的情形中它将是哈密顿空间。所以断言是：这个空间生活在理论 $\cZ$ 的真空模空间之上。也就是说，我们已经取了一个理论 $\cZ$ 并从中造出代数几何；现在我们要为一个理论 $\cZ$ 连同边界条件 $T$ 一并造代数几何。

回想我们对二维理论做过类似的事。若有一个范畴 $\cC$（我们把它看成对应二维 TFT），可以考虑一个对象 $\cF\in\cC$。然后看 $\cF$ 的自同态代数，它接收来自二维理论 $\cZ$ 的局部算子的映射：
\begin{equation}\tag{16.91}
\mathrm{HH}(\cC)=\cZ(S^1)\to\End(\cF).
\end{equation}
于是我们有一个局部算子上的代数，即它层化在局部算子的 $\Spec$ 上，奇异支撑正是由此而来。

用代数的语言重复这一点：有一个代数 $A$ 作用在模 $M$ 上，由此得到映射
\begin{equation}\tag{16.92}
\mathrm{center}(A)\to\End(M),
\end{equation}
这就是我们将要得到的那类东西。在此情形中 $\End(M)$ 只是结合（$\mathbb{E}_1$）代数，所以其实不太想对它取 $\Spec$，因为它不够交换。但在四维中一切都更交换。于是我们可以取 $\Spec$ 并看到几何。

设 $\cZ$ 是四维 TFT，$T$ 是边界条件。为热身，先假设 $\cZ$ 平凡。这种情形下 $T$ 等同于一个通常的三维 TFT。那么 $T$ 的局部算子 $\mathcal{T}(S^2)$ 构成 $\mathbb{E}_3$-代数。在上同调层面它看起来只是一个（分次的、偶的）Poisson 代数。于是可以定义这个理论 $T$ 的真空模空间为
\begin{equation}\tag{16.93}
\cM^0_T=\Spec\bigl(\mathcal{T}(S^2)\bigr),
\end{equation}
我们可以把它看成（仿射的）Poisson 簇。若取一个带 $\C^\times$ 作用的普通 Poisson 簇，则其函数环是分次的，我们可以把该分次宣布为上同调分次。也就是说，取使 $\{,\}$ 的次数为 $-2$ 的 $\C^\times$-作用。所以这是一个 $\mathbb{P}_3$-簇。

\begin{srcex}{113}
$T^*X$ 总有 $\C^\times$ 的伸缩作用。它的次数为 $2$，所以其逆——Poisson 括号——的次数为 $-2$。也就是说，我们可以把它做成 $\mathbb{P}_3$-簇。
\end{srcex}

模空间 $\cM^0_T$ 总是仿射的，因为我们只看了局部算子。改而看线算子 $\mathcal{T}(S^1)$，它是 $\mathbb{E}_2$-范畴（编织的）。所得的模空间将是一个不一定仿射的 $\mathbb{E}_3$-叠。我们正在做某种（从理论 $T$ 重构 $M$ 的）重构，但也许反过来想更容易。若从分次辛的 $M$ 出发（即带有 $\C^\times$-作用，使辛形式的权为 $2$），则可以构造一个三维 TFT，称为 \term{Rozansky--Witten 理论}{Rozansky--Witten theory}。该理论的局部算子构成
\begin{equation}\tag{16.94}
\mathcal{T}_M(S^2)=(\cO(M),\{,\}),
\end{equation}
把它看成 $\mathbb{P}_3$-代数；而线算子是
\begin{equation}\tag{16.95}
\mathcal{T}_M(S^1)=\QC(M).
\end{equation}
但线算子上的张量积其实带有编织形变，所以它实际上是 $\mathbb{E}_2$ 的。这一点真正被理解的只有 $M=T^*X$ 的情形，此时我们其实用的是 Koszul 对偶：
\begin{equation}\tag{16.96}
\QC(T^*X)\leftrightarrow\QC^!(LX),
\end{equation}
而 $\QC^!(LX)$ 带有自然的 $\mathbb{E}_2$-结构。

于是我们从 $T$ 出发，得到
\begin{equation}\tag{16.97}
\cM^1\to\cM^0,
\end{equation}
然后希望 $T$ 是 $\cM^0$ 上的 RW，更可能是 $\cM^1$ 上的 RW。

\begin{srcrem}{117}
物理学家可能把这些模空间之一的 RW 称为 $T$ 的低能有效场论。
\end{srcrem}

\par\medskip\noindent\textbf{警告 3.}\ 这种重构有时会惨败。例如考虑 Chern--Simons。它没有局部算子，所以仿射版本会是一个点。它有有意思的线算子，由这个模张量范畴给出，但若试图取它的谱，也只是一个点。这个模张量范畴没有任何导出的东西，所以它无法与 $\mathbb{E}_3$-代数对话。然而 $BG$ 由 Killing 型给出 $2$-移位辛结构，而 Chern--Simons 在道德上就是 $BG$ 的 RW，但我们从 $\cM^0$ 和 $\cM^1$ 都看不到这一点。% [Tutor Note: 原文 “the BG” 指分类叠 $BG$；“killing form” 应为首字母大写的 Killing 型。此处作者强调：Chern--Simons 的线算子范畴（模张量范畴）是半单的、没有导出结构，因此不能用局部算子/线算子的谱重构出 $BG$ 的 $2$-移位辛几何；这说明本节的重构程序并非普遍适用。]

现在设我们的体理论 $\cZ$ 非平凡，仍带边界条件 $T$。要点是：我们刚才所说的一切都有相对版本。也就是说，$T$ 的局部算子 $\mathcal{T}(S^2)$ 不再是像之前那样的裸 $\mathbb{E}_3$-代数，而是范畴 $\cZ(S^2)$ 中的 $\mathbb{E}_3$-代数。一种说法是：我们可以取体中的局部算子并把它们与边界碰撞，使 $\mathcal{T}(S^2)$ 成为 $\cZ(S^3)$ 上的代数。这类似于早先的情形：若有一个模 $A\ \circlearrowleft\ M$，则 $\End(M)$ 是 $A$ 的中心上的代数。也就是说，它是范畴 $\mathrm{center}(A)\text{-}\mathrm{mod}$ 中的（$\mathbb{E}_1$）代数对象。谈论这个范畴的代数对象是合理的，因为 $A$ 的中心是 $\mathbb{E}_2$，所以它上的模范畴本身就是 $\mathbb{E}_1$ 的。

同样的图景在三维理论中也成立：边界上的局部算子 $\mathcal{T}(S^2)$ 是 $\mathbb{E}_4$-代数 $\cZ(S^3)$ 上的 $\mathbb{E}_3$-代数，即 $\mathbb{E}_3$-范畴 $\cZ(S^3)\text{-}\mathrm{mod}$ 中的 $\mathbb{E}_3$-代数对象。现在我们处于更交换的环境，可以取 $\Spec$。这意味着我们有一个映射
\begin{equation}\tag{16.98}
\begin{array}{c}
\Spec\mathcal{T}(S^2)\\[2pt]
\downarrow\\[2pt]
\cM^0=\Spec\cZ(S^3)
\end{array}
\end{equation}
这个映射是所谓拉格朗日映射的一个例子（它与 Poisson 结构有某种相容性）。范畴 $\cZ(S^3)\text{-}\mathrm{mod}$ 其实是对 $\cZ(S^2)$ 的近似。更好的想法是：$\mathcal{T}(S^2)$ 是 $\mathbb{E}_3$-范畴 $\cZ(S^2)$ 中的 $\mathbb{E}_3$-代数对象。

若身处 Langlands 语境，即体理论 $\cZ$ 实际上是 $\mathcal{B}_{G^\vee}$，我们看到：对 $\mathcal{B}_{G^\vee}$ 的边界理论 $T$，我们得到 $\fg^{\vee*}[2]/G^\vee$ 上的代数，因为这个范畴 $\cZ(S^2)$ 是余伴随表示 $\fg^{\vee*}[2]/G^\vee$ 上的层。于是取 $\Spec$ 时，得到一个仿射态射
\begin{equation}\tag{16.99}
M/G^\vee\to\fg^{\vee*}[2]/G^\vee,
\end{equation}
它与 Poisson 结构相容。这等同于给出一个仿射哈密顿 $G^\vee$-簇（在分次剪切意义下）$G^\vee\ \circlearrowleft\ M$。然后若想看到非仿射的东西，可以看 $\mathcal{T}(S^1)$，它是 $\cZ(S^2)$ 上的 $\mathbb{E}_2$-代数。于是我们得到类似
\begin{equation}\tag{16.100}
\QC(M/G^\vee)
\end{equation}
的东西，它被
\begin{equation}\tag{16.101}
\QC(\fg^{\vee*}[2]/G^\vee)
\end{equation}
通过矩映射作用。所以 $\mathcal{B}_{G^\vee}$ 的边界条件以哈密顿 $G^\vee$-空间作为其“影子”。反过来，若 $G^\vee\ \circlearrowleft\ M$ 是全纯辛簇上的哈密顿作用，则 $\mathrm{RW}_M$ 定义了 $\mathcal{B}_{G^\vee}$ 的一个边界条件。
